%%% Hlavní soubor. Zde se definují základní parametry a odkazuje se na ostatní části. %%%

% Meta-data o práci (je nutno upravit)
\input metadata.tex

% Vygenerujeme metadata ve formátu XMP pro použití balíčkem pdfx
\input xmp.tex

%% Verze pro jednostranný tisk:
% Okraje: levý 40mm, pravý 25mm, horní a dolní 25mm
% (ale pozor, LaTeX si sám přidává 1in)
\documentclass[12pt,a4paper]{report}
\setlength\textwidth{145mm}
\setlength\textheight{247mm}
\setlength\oddsidemargin{15mm}
\setlength\evensidemargin{15mm}
\setlength\topmargin{0mm}
\setlength\headsep{0mm}
\setlength\headheight{0mm}
% \openright zařídí, aby následující text začínal na pravé straně knihy
\let\openright=\clearpage

%% Pokud tiskneme oboustranně:
% \documentclass[12pt,a4paper,twoside,openright]{report}
% \setlength\textwidth{145mm}
% \setlength\textheight{247mm}
% \setlength\oddsidemargin{14.2mm}
% \setlength\evensidemargin{0mm}
% \setlength\topmargin{0mm}
% \setlength\headsep{0mm}
% \setlength\headheight{0mm}
% \let\openright=\cleardoublepage

%% Pokud práci odevzdáváme pouze elektronicky, vypadají lépe symetrické okraje
% \documentclass[12pt,a4paper]{report}
% \setlength\textwidth{145mm}
% \setlength\textheight{247mm}
% \setlength\oddsidemargin{10mm}
% \setlength\evensidemargin{10mm}
% \setlength\topmargin{0mm}
% \setlength\headsep{0mm}
% \setlength\headheight{0mm}
% \let\openright=\clearpage

%% Vytváříme PDF/A-2u
\usepackage[a-2u]{pdfx}

%% Přepneme na českou sazbu a fonty Latin Modern
\usepackage[czech]{babel}
\usepackage{lmodern}

% Pokud nepouživáme LuaTeX, je potřeba ještě nastavit kódování znaků
\usepackage{iftex}
\ifpdftex
\usepackage[utf8]{inputenc}
\usepackage[T1]{fontenc}
\usepackage{textcomp}
\fi

%%% Další užitečné balíčky (jsou součástí běžných distribucí LaTeXu)
\usepackage{amsmath}        % rozšíření pro sazbu matematiky
\usepackage{amsfonts}       % matematické fonty
\usepackage{amsthm}         % sazba vět, definic apod.
\usepackage{bm}             % tučné symboly (příkaz \bm)
\usepackage{booktabs}       % lepší vodorovné linky v tabulkách
\usepackage{caption}        % umožní definovat vlastní popisky plovoucích objektů
\usepackage{csquotes}       % uvozovky závislé na jazyku
\usepackage{dcolumn}        % vylepšené zarovnání sloupců tabulek
\usepackage{floatrow}       % umožní definovat vlastní typy plovoucích objektů
\usepackage{graphicx}       % vkládání obrázků
\usepackage{icomma}         % inteligetní čárka v matematickém módu
\usepackage{indentfirst}    % zavede odsazení 1. odstavce kapitoly
\usepackage[nopatch=item]{microtype}  % mikrotypografická rozšíření
\usepackage{paralist}       % lepší enumerate a itemize
\usepackage[nottoc]{tocbibind} % zajistí přidání seznamu literatury,
                            % obrázků a tabulek do obsahu
\usepackage{xcolor}         % barevná sazba

% Balíček hyperref, kterým jdou vyrábět klikací odkazy v PDF,
% ale hlavně ho používáme k uložení metadat do PDF (včetně obsahu).
% Většinu nastavítek přednastaví balíček pdfx.
\hypersetup{unicode}
\hypersetup{breaklinks=true}

% Balíčky pro sazbu informatických prací
\usepackage{algpseudocode}  % součást balíčku algorithmicx
\usepackage[Algoritmus]{algorithm}
\usepackage{fancyvrb}       % vylepšené prostředí verbatim
\usepackage{listings}       % zvýrazňování syntaxe zdrojových textů

% Cleveref může zjednodušit odkazování, ale jeho užitečnost pro češtinu
% je minimalní, protože nezvládá skloňování.
% \usepackage{cleveref}

% Formátování bibliografie (odkazů na literaturu)
% Detailní nastavení můžete upravit v souboru macros.tex.
%
% POZOR: Zvyklosti různých oborů a kateder se liší. Konzultujte se svým
% vedoucím, jaký formát citací je pro vaši práci vhodný!
%
% Základní formát podle normy ISO 690 s číslovanými odkazy
\usepackage[natbib,style=iso-numeric,sorting=none]{biblatex}
% ISO 690 s alfanumerickými odkazy (zkratky jmen autorů)
%\usepackage[natbib,style=iso-alphabetic]{biblatex}
% ISO 690 s citacemi tvaru Autor (rok)
%\usepackage[natbib,style=iso-authoryear]{biblatex}
%
% V některých oborech je běžnější obyčejný formát s číslovanými odkazy
% (sorting=none říká, že se bibliografie má řadit podle pořadí citací):
%\usepackage[natbib,style=numeric,sorting=none]{biblatex}
% Číslované odkazy, navíc se [1,2,3,4,5] komprimuje na [1-5]
%\usepackage[natbib,style=numeric-comp,sorting=none]{biblatex}
% Obyčejný formát s alfanumerickými odkazy:
%\usepackage[natbib,style=alphabetic]{biblatex}

% Custom přidané packages a definice
\usepackage{tikz}
\usetikzlibrary{shapes.geometric, arrows}
\usepackage{svg}
\usepackage{float}
%\newcommand{\todo}[1]{\textcolor{red}{#1}}
\DeclareRobustCommand{\todo}[1]{\textcolor{red}{#1}}

\DeclareRobustCommand{\notSureIfOkToSay}[1]{\textcolor{orange}{#1}} % --> inaccurate or non well formulated hypothesis

\DeclareRobustCommand{\notSureIfOkToSayThisWay}[1]{\textcolor{blue}{#1}} % --> formal, grammar or syntactic thing in text ...

\usepackage{siunitx}

\sisetup{
    output-decimal-marker = {,},
    detect-weight = true,
    detect-family = true,
    round-mode = places,
    round-precision = 4,
    group-digits = false,
    mode = math
}

\newcommand{\bnum}[1]{\bfseries\tablenum{#1}}

%\usepackage{tcolorbox}
%\newtcolorbox{todobox}{colback=red!5!white, colframe=red!75!black}

% Z tohoto souboru se načítají položky bibliografie
\addbibresource{literatura.bib}

% Definice různých užitečných maker (viz popis uvnitř souboru)
\input macros.tex

%%% Titulní strana a různé povinné informační strany
\begin{document}
\include{title}

%%% Strana s automaticky generovaným obsahem práce

\tableofcontents

%%% Jednotlivé kapitoly práce jsou pro přehlednost uloženy v samostatných souborech
\include{uvod}
\include{kap01}
\include{kap02}
\include{kap03}
\include{kap04}
\include{kap05}
\include{kap06}
\include{kap07}
\include{kap08}

\include{zaver}

%%% Seznam použité literatury
\include{literatura}

%%% Obrázky v práci
%%% (pokud jich je malé množství, obvykle není třeba seznam uvádět)
\listoffigures

%%% Tabulky v práci (opět nemusí být nutné uvádět)
%%% U matematických prací může být lepší přemístit seznam tabulek na začátek práce.
\listoftables

%%% Použité zkratky v práci (opět nemusí být nutné uvádět)
%%% U matematických prací může být lepší přemístit seznam zkratek na začátek práce.
\chapwithtoc{Seznam použitých zkratek}

%%% Součástí doktorských prací musí být seznam vlastních publikací
\ifx\ThesisType\TypePhD
\chapwithtoc{Seznam publikací}
\fi

%%% Přílohy k práci, existují-li. Každá příloha musí být alespoň jednou
%%% odkazována z vlastního textu práce. Přílohy se číslují.
%%%
%%% Do tištěné verze se spíše hodí přílohy, které lze číst a prohlížet (dodatečné
%%% tabulky a grafy, různé textové doplňky, ukázky výstupů z počítačových programů,
%%% apod.). Do elektronické verze se hodí přílohy, které budou spíše používány
%%% v elektronické podobě než čteny (zdrojové kódy programů, datové soubory,
%%% interaktivní grafy apod.). Elektronické přílohy se nahrávají do SISu.
%%% Povolené formáty souborů specifikuje opatření rektora č. 72/2017.
%%% Výjimky schvaluje fakultní koordinátor pro zavěrečné práce.
\appendix
\chapter{Přílohy}

\section{Přehled základních agregačních funkcí}


\begin{table}[b!]
    \centering
    \begin{tabular}{l@{\hspace{1.3cm}}D{.}{,}{3.2}D{.}{,}{2.3}D{.}{,}{2.3}D{.}{,}{2.3}D{.}{,}{2.3}D{.}{,}{2.3}}
    \toprule
    \textbf{Člen} &
    \mc{\textbf{$i_1$}} &
    \mc{\textbf{$i_2$}} &
    \mc{\textbf{$i_3$}} &
    \mc{\textbf{$i_4$}} &
    \mc{\textbf{$i_5$}} &
    \mc{\textbf{$i_6$}} \\
    \midrule
    Alice & 10 & 8 & 4 & 4 & 5 & 7 \\
    Bob & 7 & 2 & 5 & 4 & 4 & 3 \\
    Carol & 6 & 1 & 8 & 4 & 1 & 2 \\
    \bottomrule
    \end{tabular}
    \caption{Výchozí tabulka hodnocení předmětů katalogu $i_1-i_6$}\label{tab:ScoresForApendix}
\end{table}

\begin{table}[b!]
    \centering
    \begin{tabular}{| p{3cm} | p{5cm} | p{5cm} | p{5cm} |}
    \toprule
    \textbf{Agregační strategie} &
    \mc{\textbf{Jak funguje}} &
    %\mc{\textbf{Kategorie}} &
    \mc{\textbf{Příklad}} \\
    \midrule
    \midrule \textit{Plurality voting/majority voting}
    & Opakovaně aplikuje \textit{first past the post}, položka s nejvíce hlasy je vybrána.
    & Vyhrává položka $i_1$, protože majorita skupiny položce dala nejvyšší hodnocení.
    Dále již bez $i_1$ vyhraje $i_3$ atd. \\
    %\midrule Majority voting & Vyhrává položka s nadpoloviční většinou hlasů. & První vybraná položka je item $i_1$ - protože má nejvyšší počet hlasů, v dalším kole bychom vybrali položku $i_2$ a takto pokračujeme. \\
    \midrule \textit{Average} & Z průměrujeme jednotlivá hodnocení všech členů. & Skupinové hodnocení pro $i_3$~je $(4+5+8)/3 = 6$. \\
    \midrule \textit{Multiplicative} & Jednotlivá hodnocení vynásobíme. & Pro $i_3$ je skóre $4\cdot5\cdot8 = 160$.  \\
    \midrule \textit{Borda count} & Položkám přiřadíme skóre podle relevance pro každého uživatele.
    Následně sečteme jednotlivá skóre přes všechny uživatele.  &

    %Pro uživatele A dostaneme řazení dle skóre\footnote{
    %    Položce udělíme skóre odpovídající pozici v seznamu sestupně. Tedy pro $n$ položek
    %    první předmět s nejvyšším hodnocením dostane $n-1$, druhý $n-2$, ...
    %    Pokud některé položky mají stejné hodnocení, tak jejich skóre zprůměrujeme.
    %}.
    Pro Alici dostaneme řazení podle skóre:
    (($i_1$, 5), ($i_2$, 4), ($i_6$, 3), ($i_5$, 2), ($i_3$, 0.5), ($i_4$, 0.5)),
    zbylá skóre jsou pro přehlednost v \ref{tab:BordaScore}.
    Pro skupinu dostaneme výsledné řazení: (($i_1$, 14), ($i_3$, 9.5), ($i_4$, 6), ($i_5$, 5), ($i_6$, 5), ($i_2$, 5)).\\
    %(($i_1$, 14), ($i_3$, 9.5), ($i_4$, 6), ($i_5$, 5), ($i_6$, 5), ($i_2$, 5))
    \midrule \textit{Copeland rule} &
    Spočítá kolikrát předmět v majoritním hlasování porazí
    ostatní položky od čehož odečítá kolikrát prohraje.
    & Položka $i_3$ má skupinové skóre 3,
    protože jeho majoritní hodnocení je 17,
    což je vyšší než vychází pro skupinky
    ($i_2$, $i_4$, $i_5$, $i_6$), naopak ($i_1$) má hlasování vyšší, tedy 4 - 1 = 3 \\
    \midrule \textit{Approval voting} & Počítá počet individuálních hlasů vyšších než je stanovená hranice. & Pro hranici $2$ a položku $i_3$ dostaneme $3$ a pro $i_5$ je $2$.\\
    \midrule \textit{Least Misery} & Vybírá minimum z hlasů. & $i_1$ tedy dostane $6$.  \\
    \midrule \textit{Most pleasure} & Vybírá maximum z hlasů. & Pro $i_1$ je to $10$.  \\
    \midrule \textit{Average without Misery} & Průměr, ale vynecháme hlasy nižší než je stanovená hranice. &
    Pro min. hranici 2 vynecháme $i_2$, $i_5$, protože někteří uživatelé je nesnáší. \\
    \midrule \textit{Fairness} & Postupně vybíráme z maximálních hodnocení jednotlivých členů. &
    Vyhrává $i_1$ (nejlépe hodnocené Alicí), potom $i_3$ (nejlepší od Boba) a na konec
    $i_6$ (další nejlepší od Carol) a opakujeme.  \\
    \midrule
    \textit{Most respected person} & Používá hodnocení nejrespektovanější osoby ve skupině. &
    Pokud je Carol nejrespektovanější osoba, potom $i_1$ má hodnocení 6.\\
    \midrule
    \bottomrule
    \end{tabular}

    \caption{Shrnutí možných agregačních strategií vycházející z teorie sociálního výběru.}
    \label{tab:AllBasicStrategies}
\end{table}


\begin{table}[b!]
    \centering
    \begin{tabular}{l@{\hspace{1.3cm}}D{.}{,}{3.2}D{.}{,}{2.3}D{.}{,}{2.3}D{.}{,}{2.3}D{.}{,}{2.3}D{.}{,}{2.3}}
    \toprule
    \textbf{Člen} &
    \mc{\textbf{$i_1$}} &
    \mc{\textbf{$i_2$}} &
    \mc{\textbf{$i_3$}} &
    \mc{\textbf{$i_4$}} &
    \mc{\textbf{$i_5$}} &
    \mc{\textbf{$i_6$}} \\
    \midrule
    Skóre pro Alici & 5 & 4 & 0.5 & 0.5 & 2 & 3 \\
    Skóre pro Boba   & 5 & 0 & 4 & 2.5 & 2.5 & 1 \\
    Skóre pro Carol & 4 & 0.5 & 5 & 3 & 0.5 & 2 \\
    Skóre skupiny & 14 & 4.5 & 9.5 & 6 & 5 & 5 \\
    \bottomrule
    \end{tabular}
    \caption{Borda count skóre pro položky $i_1 - i_6$. Pokud má více prvků stejné
    hodnocení, každý dostane průměrné skóre. }
    \label{tab:BordaScore}
\end{table}



\begin{table}[b!]
    \centering
    \begin{tabular}{l@{\hspace{1.5cm}}l}
    \toprule
    \textbf{Typy restaurací}\\
    \midrule
    restaurant & spanish\_restaurant \\
    donut\_shop & turkish\_restaurant \\
    bakery & wine\_bardeli \\
    fast\_food\_restaurant & vietnamese\_restaurant \\
    breakfast\_restaurant & ramen\_restaurant \\
    chocolate\_factory & steak\_house \\
    hamburger\_restaurant & cafe \\
    brunch\_restaurant & buffet\_restaurant \\
    dog\_cafe & fine\_dining\_restaurant \\
    japanese\_restaurant & cat\_cafe \\
    dessert\_restaurant & korean\_restaurant \\
    juice\_shop & deli \\
    diner & tea\_house \\
    sushi\_restaurant & brazilian\_restaurant \\
    italian\_restaurant & indian\_restaurant \\
    pub & mexican\_restaurant \\
    meal\_delivery & pizza\_restaurant \\
    bar\_and\_grill & afghani\_restaurant \\
    chinese\_restaurant & middle\_eastern\_restaurant \\
    vegan\_restaurant & greek\_restaurant \\
    mediterranean\_restaurant & ice\_cream\_shop \\
    dessert\_shop & cafeteria \\
    wine\_bar & food\_court \\
    american\_restaurant & french\_restaurant \\
    vegetarian\_restaurant & seafood\_restaurant \\
    thai\_restaurant &  \\
    chocolate\_shop &  \\
    barbecue\_restaurant &  \\
    bagel\_shop &  \\
    asian\_restaurant &  \\
    bar &  \\
    confectionery &  \\
    indonesian\_restaurant &  \\
    lebanese\_restaurant &  \\
    sandwich\_shop &  \\
    african\_restaurant &  \\
    \bottomrule
    \end{tabular}
    \caption{Klíčová slova vybraná pro vyhledávání }
    \label{tab:FoodEstablishmentTypes}
\end{table}


\begin{listing}[H]
    \begin{lstlisting}
    {
        id - id mista,
        displayName - nazev mista,
        formattedAddress - adresa mista,
        location - GPS souradnice,
        types - typy prirazene dle ciselniku,
        photos - seznam jmen fotek mista,
        rating - hodnoceni mista z google maps,
        reviews - az 3 uzivatelska reviews,
        priceLevel - cenova hladina (dle ciselniku),
        allowsDogs - zdali povoluje vstup pro psi,
        googleMapsUri - odkaz na Google maps s detailem mista,
        regularOpeningHours - bezna oteviraci doba,
        goodForGroups - je vhodne pro skupiny,
        goodForChildren - je vhodne pro deti,
        outdoorSeating - zdali nabizi venkovni sezeni,
        delivery - zdali nabizi doruceni,
        takeout - zdali nabizi jidlo sebou,
        liveMusic - zdali nabizi zivou hudbu,
        servesVegetarianFood - zdali nabizi vegetarianske jidlo
    }
\end{lstlisting}
\caption{Struktura \textit{body} dotazu pro \textit{Places text search (New)} API.}
\label{prog3:usedFieldMaskFields}
\end{listing}



\clearpage

\section{Analýza samplovacího algoritmu inicializačních dat}
\label{initSamplingAlgorithmAnalysis}
V této části provedeme analýzu různých metod
samplování filmů z \textit{datasetu MovieLens} pro inicializační session.

\subsection{Cluster sampling (content based)}
Jednou z možností jak zajistit
dobrou diverzifikovanost vzorků
je data nejprve rozdělit do \textit{clusterů},
nad kterými následně provedeme výběr.

Nejprve je však nutné vytvořit vhodné
reprezentace filmů. Ty můžeme zkonstruovat
\textit{one-hot-encodingem} příslušných
žánrů\footnote{Mohlo by nás také napadnout použít MovieLens tagy,
nicméně těch je velké množství, což by vedlo k
vysoké dimenzionalitě a problémům při \textit{clustrování}.}.
Pro obohacení navíc můžeme přidat
dimenze popularity a průměrného hodnocení
(škálované na [0,1]).
Film~$i$~má~tedy reprezentaci

\begin{equation}
    r(i) = [genres_i, popularity_i, average\_rating_i].
\end{equation}

Následně provedeme rozdělení do \textit{k} \textit{clusterů} metodou
\textit{k-means}. Pro určení optimální hodnoty \textit{k} se lze podívat na dvě metriky.
Tou první je určení součtu čtverců vnitřních \textit{clusterů} (WCSS)
pro různých velikosti \textit{k} a pokusit se na grafu nalézt
tzv. \textit{elbow}\footnote{Elbow - bod největšího zakřivení průběhu.}.

%\begin{equation}
%    \text{WCSS} = \sum_{i=1}^{k} \sum_{x \in C_i} \| x - \mu_i \|^2
%\end{equation}

\begin{figure}[H]\centering
    \includegraphics[width=145mm]{img/clustering-kmeans-elbow-method.pdf}
    \caption{
        Hodnoty WCSS pro různé hodnoty \textit{k} nad datasetem MovieLens.
        }
    \label{obr20:elbowMethod}
\end{figure}


Z výsledků analýzy metody \textit{elbow} -- graf \ref{obr20:elbowMethod}
lze vypozorovat, že v grafu nedostáváme ostrý loket
a máme tak pouze hrubý odhad $ k \in (18, 25) $.
To naznačuje, že se \textit{clustry}
do značné míry v \textit{datasetu} překrývají
a nelze je tak jednoznačně určit.

Pro konkrétní hodnotu můžeme použít \textit{KneeLocator}
z python balíčku \textit{kneed}, který v křivce
identifikuje největší změnu zakřivení.
Odtud dostaneme hodnotu~\mbox{$k=20$}.

\begin{figure}[H]\centering
    \includegraphics[width=145mm]{img/clustering-silueth-score.pdf}
    \caption{
        Hodnoty WCSS pro různé hodnoty $k$ nad datasetem MovieLens.
        }
    \label{obr21:siluethScore}
\end{figure}


Alternativně lze pro různé hodnoty $k$
změřit průměrný \textit{silueth-score}, který určuje
jak dobře každý bod do svého \textit{clusteru} náleží.
Výsledný průběh \ref{obr21:siluethScore} ukazuje, že
významný růst se zastaví při hodnotě $k=20$,
kdy křivka osciluje zhruba okolo hodnoty~$4,2$.



Protože se data v \textit{datasetu} překrývají
a některé filmy tak mohou náležet do více \textit{clustrů} současně,
mohli bychom zkusit použít metodu \textit{soft} \textit{clustrování}.
Konkrétně můžeme provést \textit{fuzzy-kmeans}.
Hodnotu $k$ ponecháme, parametr \textit{fuzzy} stanovíme
standardně na 2. To nám pro každý film určí pravděpodobnostní
distribuci \textit{příslušnosti k jednotlivým clusterům}.

Následně budeme samplovat 24 filmů přes jednotlivé clustery
metodou \textit{round-robin}.

\subsection{Cluster sampling (user based)}
Alternativně bychom mohli \textit{clustrovat}
uživatele podle jejich profilů hodnocení.
To nám umožní rozdělit uživatele
do skupin dle skrytých preferencí, které
nedokážeme z \textit{datasetu} jednoduše
určit (a tedy je nelze jednoduše
zakódovat do reprezentace filmů,
jako například žánry).

Specificky pro tento případ,
můžeme původní \textit{clustrovací}
algoritmus \textit{kmeans} upravit
tak, aby jako metriku pro přiřazení
bodů do \textit{clusterů} nepoužíval \textit{euklidovskou metriku},
%tedy, že uživatel $u \in C_i \quad \text{where} \quad i = \arg\min_j \|u - \mu_j\|^2$,
ale zvolit \textit{cluster} podle \textit{Pearsonovy similarity},
která je pro určení podobnosti uživatel vhodnější.

Samplování poté můžeme provést tak, že doporučíme dostatečně
velký vzorek kandidátů (například 100), ze kterého
samplujeme uniformě náhodně nejpopulárnější
předměty.







\subsection{Baseline predictor}
Alternativně můžeme nad daty použít některý z
\textit{model based} doporučovacích přístupů,
pomocí kterého vygenerujeme dostatečný počet kandidátů
z nichž budeme samplovat. Protože chceme generovat
doporučení na základě celého datasetu, tak se jako vhodný
kandidát jeví \textit{base-line} prediktor
$$
    \hat{r}_{ui} = \mu + b_u + b_i,
$$

kde $b_u$ -- \textit{user bias}, vyjadřující okolik se uživatelské hodnocení liší od globální průměru
a $b_i$ -- \textit{item bias}, vyjadřující okolik se hodnocení předmětu liší od globálního průměru.
Tyto složky jsou určeny vyřešením následující optimalizační funkce

$$
\min_{b_u, b_i} \sum_{(u,i) \in \mathcal{K}} (r_{ui} - \mu - b_u - b_i)^2 + \lambda(b_u^2 + b_i^2),
$$
kde $\mathcal{K}$ je množina hodnocených předmětů uživatelem $u \in \mathcal{U}$ a $\lambda$ -- parametr L2 regularizace.

\subsection{Metriky evaluace}


\begin{figure}[H]\centering
    \includegraphics[width=145mm]{img/init-sampling-comparison.pdf}
    \caption{
        Porovnání jednotlivých metod na stanovených metrikách.
        }
    \label{obr21:initAlgoComparisionApendix}
\end{figure}





\section{Parametry upravené sigmoid funkce}
\label{sec:ShiftedSigmoid}

Upravená \textit{sigmoid} funkce $\sigma_{c_0}: \mathbb{R} \rightarrow (c_0, 1)$ s dolním omezením $c_0 \in [0, 1)$,
posunem $k_0 \in \mathbb{R}$ a parametrem $\alpha \in \mathbb{R}$ v $k \in \mathbb{R}$ má tvar

\begin{equation}
    \sigma_{c_0}(k) = \left(c_0 + (1-c_0) \frac{1}{1+e^{-\alpha(k-k_0)}} \right).
\end{equation}


Označme $\sigma' = \sigma_{c_0}(k)$. Jednoduchou úpravou pro $k$ dostaneme

%Převedeme~výrazy~s~$c_0$~na pravou stranu a inverzí $()^{-1}$ obou stran dostaneme
%$$
%\frac{1-c_0}{\psi - c_0} = 1 + e^{-\alpha(k-k_0)},
%$$
%po odečtení 1 a aplikací logaritmu $\ln()$ na obě strany
%$$
%\ln\left(\frac{1-\psi}{\psi - c_0}\right) =  -\alpha(k-k_0),
%$$
%zjevně tedy
$$
k = k_0 - \frac{1}{\alpha} \ln\left( \frac{1-\sigma'}{\sigma'-c_0} \right).
$$

%$$
%\psi(k) = \left(c_0 + (1-c_0) \frac{1}{1+e^{-\alpha(k-k_0)}} \right)
%$$

Pro $\sigma' = c_0 + \epsilon$ a nějaké malé $\epsilon \in (0, 1-c_0)$ dostaneme
\begin{equation}
    k = k_0 - \delta_{\epsilon},
\end{equation}
kde
\begin{equation}
    \delta_{\epsilon} = \frac{1}{\alpha} \ln\left( \frac{1-c_0}{\epsilon} - 1 \right).
\end{equation}

Horizontální posun funkce $\sigma'$ se zhruba konstantním růstem $c_0$ s přesností $\epsilon$ můžeme
tedy označit $k_{\epsilon} := k_0 + \delta_{\epsilon}$.

\medbreak

\subsubsection{Další užitečná pozorování:}
\begin{itemize}
    \item \textit{Inflexní} bod $\sigma_{c_0}$ má souřadnice $((1+c_0)/2, k_0/\alpha)$.
\end{itemize}


\subsection{Horní mez upravené sigmoid funkce}
\label{sec:ShiftedSigmoidUpper}



Pro horní omezení \textit{sigmoid} funkce $\sigma_{c_0}$ můžeme buď upravit výraz $(1-c_0)$
na obecnější tvar $(c-c_0)$, kde $c: \mathbb{N}_0 \rightarrow [c_0, 1]$,
nebo si lze uvědomit, že stačí vhodně přeškálovat část se \textit{sigmoid} funkcí,
tedy stačí nalézt vhodné $\gamma: \mathbb{N}_0 \rightarrow [0, 1]$,
protože potom dostaneme
\begin{equation}
    \sigma'_{c_0}(k, s) = c_0 + (1-c_0)\sigma(k)\gamma(s),
\end{equation}
kde $\sigma$ je obecný \textit{sigmoid} v $k \in \mathbb{R}$ a $s \in \mathbb{N}_0$ je horní omezení.
Pro účely práce uvažujme $s \in [s_{min}, s_{max}] \subset \mathbb{N}_0$.

%$$
%    f(x) = c_0 + (1-c_0)f(x) - (1-c_0)f(x)\gamma,
%$$
%tedy, že:
%$$
%    f(x) = c_0 + (1-c_0)f(x)(1 - \gamma).
%$$
%Neboli pro horní mez stačí nalézt vhodnou funkci $\gamma: \mathbb{R} \rightarrow [0, 1]$

\subsubsection{Vhodné volby $\gamma$:}
\begin{equation}
    \gamma(s) = \left(\frac{\max\left(0,\min\left(s,s_{max}\right)-s_{min}\right)} {s_{max}-s_{min}}\right)^p
\end{equation}
pro $p\in \mathbb{N}$

\begin{equation}
    \gamma(s) = \frac{ e^{ \beta\max\left(0,\min\left(s,s_{max}\right)-s_{min}\right)} - 1 }{e^{ \beta\left( s_{max}-s_{min} \right)} - 1}
\end{equation}
pro $\beta \in \mathbb{R^+}$

\clearpage

\begin{table}[b!]
    \centering
    \begin{tabular}{l@{\hspace{1.5cm}}D{.}{,}{1.5}@{\hspace{1.5cm}}D{.}{,}{1.3}@{\hspace{1.5cm}}D{.}{,}{1.5}}
    \toprule
     & \mc{} & \mc{} & \mc{} \\
    \pulrad{\textbf{$\lambda$}} & \mc{\pulrad{\textbf{precision@20}}} & \mc{\pulrad{\textbf{recall@20}}} & \mc{\pulrad{\textbf{NDCG@20}}} \\
    \midrule
    100 & 0.00132 & 0.024 & 0.026 \\
    200 & 0.00137 & 0.025 & 0.027 \\
    300 & 0.00142 & 0.027 & 0.028 \\
    400 & 0.00157 & 0.029 & 0.031 \\
    500 & 0.00163 & 0.03 & 0.033 \\
    600 & 0.00158 & 0.029 & 0.032 \\
    700 & 0.00148 & 0.027 & 0.03 \\
    800 & 0.00142 & 0.026 & 0.028 \\
    900 & 0.00147 & 0.026 & 0.029 \\
    1000 & 0.00142 & 0.025 & 0.028 \\
    1100 & 0.00142 & 0.025 & 0.028 \\
    1200 & 0.00153 & 0.028 & 0.031 \\
    1300 & 0.00142 & 0.025 & 0.028 \\
    1400 & 0.00142 & 0.025 & 0.028 \\
    1500 & 0.00147 & 0.026 & 0.029 \\
    1600 & 0.00142 & 0.025 & 0.028 \\
    1700 & 0.00147 & 0.026 & 0.029 \\
    1800 & 0.00153 & 0.028 & 0.031 \\
    1900 & 0.00142 & 0.025 & 0.028 \\
    \bottomrule
    \end{tabular}
    \caption{Přesnost $\text{EASE}^R$ pro různé volby $\lambda$ pro \textit{dataset} s minimálním počtem hodnocení pro uživatele $\#(r_{min}) \ge 2$}
    \label{tab:EaserRegMeasure}
    \end{table}

    \begin{table}[b!]
    \centering
    \begin{tabular}{l@{\hspace{1.5cm}}D{.}{,}{1.5}@{\hspace{1.5cm}}D{.}{,}{1.3}@{\hspace{1.5cm}}D{.}{,}{1.5}}
    \toprule
     & \mc{} & \mc{} & \mc{} \\
    \pulrad{\textbf{$\lambda$}} & \mc{\pulrad{\textbf{precision@20}}} & \mc{\pulrad{\textbf{recall@20}}} & \mc{\pulrad{\textbf{NDCG@20}}} \\
    \midrule
    100 & 0.00163 & 0.03 & 0.033 \\
    200 & 0.00172 & 0.031 & 0.034 \\
    300 & 0.00172 & 0.031 & 0.034 \\
    400 & 0.00172 & 0.031 & 0.034 \\
    500 & 0.0018 & 0.033 & 0.036 \\
    600 & 0.00172 & 0.031 & 0.034 \\
    700 & 0.0018 & 0.033 & 0.036 \\
    800 & 0.0018 & 0.033 & 0.036 \\
    900 & 0.00172 & 0.031 & 0.034 \\
    1000 & 0.00181 & 0.032 & 0.036 \\
    1100 & 0.00181 & 0.032 & 0.036 \\
    1200 & 0.00172 & 0.031 & 0.034 \\
    1300 & 0.0018 & 0.033 & 0.036 \\
    1400 & 0.00189 & 0.034 & 0.038 \\
    1500 & 0.00172 & 0.031 & 0.034 \\
    1600 & 0.00181 & 0.032 & 0.036 \\
    1700 & 0.0018 & 0.033 & 0.036 \\
    1800 & 0.00181 & 0.032 & 0.036 \\
    1900 & 0.00181 & 0.032 & 0.036 \\
    \bottomrule
    \end{tabular}
    \caption{Přesnost $\text{EASE}^R$ pro různé volby $\lambda$ pro \textit{dataset} s minimálním počtem hodnocení pro uživatele $\#(r_{min}) \ge 3$}
    \label{tab:EaserRegMeasure}
\end{table}

%Pro horní omezení \textit{sigmoid} funkce $\sigma_{c_0}$ můžeme upravit výraz $(1-c_0)$
%na obecnější tvar $(c-c_0)$, kde $c: \mathbb{N}_0 \rightarrow [c_0, 1]$.
%
%Mějme $s \in [s_{min}, s_{max}] \subseteq \mathbb{N}_0$ a funkci $\gamma: \mathbb{N}_0 \rightarrow [0, 1]$,
%potom lze $c$ obecně vyjádřit jako:
%
%\begin{equation}
%    c(s) = \min \left({1, \max\left({c_0, \gamma(s)}\right)} \right).
%\end{equation}
%
%\textbf{Možné volby $\gamma$:}
%\begin{equation}
%    \gamma(s) = \left( \frac{\min\left(s, s_{max}\right)-s_{min}}{s_{max}-s_{min}} + c_0 \right)^p,
%\end{equation}
%pro $p\in \mathbb{N}_0$
%
%\begin{equation}
%    \gamma(s) = \frac{ e^{ \left( \min\left(s, s_{max}\right)-s_{min} \right)} - 1 }{e^{ \left( s_{max}-s_{min} \right)} - 1} + c_0
%\end{equation}
%

%Alternativně lze funkci $\sigma_{c_0}$ vhodně přeškálovat
%\begin{equation}
%    \gamma(s) = \max\left(0, \frac{\min\left(s, s_{max}\right)-s_{min}}{s_{max}-s_{min}}\right)
%\end{equation}




%Uživaleli chceme nabídnout ideálně co možná nejširší spektrum
%\section{Vývoj preferencí v čase}
%Další problematiku, kterou jsme zde nerozebrali, ale je také dobré
%se nad ni zamyslet je fakt, že uživatelské preference se v čase vyvíjejí.
%Může nás například napadnou otázka, co když si uživatel v naší aplikaci
%již vytvoří nějaký dobře odpovídající profil a poté bude delší dobu,
%třeba i několik let, neaktivní. V takovém případě by mohlo dávat
%smysl celou inicializační fázi zopakovat. Původní profil bychom
%poté buď kompletně nahradili, nebo bychom mohli provést nejprve
%podobnostní analýzu a například cosinouvou podobností určit,
%jak moc se nově inicializovaná data vůči historickým datům
%liší a na základě toho systém aktualizovat.
%



\end{document}
